\section{Introduction}
\label{sec:intro}

Microservice architectures deployed on Kubernetes are in permanent flux. Rolling deployments,
horizontal autoscaling, configuration changes and traffic ramps all shift the statistical
distribution of observability signals without any underlying fault. Anomaly detectors that react
to distributional change therefore fire on these \emph{benign drifts} exactly as they fire on
real incidents. We confirm this directly: a per-feature $z$-score detector raises a false alarm
on $100\%$ of benign-drift episodes at every threshold $\sigma\in[2.0,3.5]$ we tested
(Section~\ref{sec:results}). The conflation of drift and anomaly is the operational problem EWAT
attacks.

\paragraph{Early warning, not RCA.}
EWAT is an \emph{early-warning} system: it predicts \emph{what} type of fault is developing and
\emph{how long before} it manifests. It is not root-cause analysis (RCA). RCA is post-mortem and
answers \emph{where} and \emph{why} a failure occurred; EWAT acts \emph{before} the failure and
answers \emph{what} and \emph{when}. Table~\ref{tab:ewat-vs-rca} makes the distinction explicit.
This boundary is a design commitment, not a slogan: the typing and ontology stages approach the
``what'' but never attribute a root cause, and our evaluation metrics (lead time, false-alarm
rate on drifts) are warning metrics, not diagnostic ones.

\begin{table}[t]
\centering
\caption{EWAT (early warning) versus root-cause analysis.}
\label{tab:ewat-vs-rca}
\small
\begin{tabularx}{\columnwidth}{@{}lXX@{}}
\toprule
 & \textbf{EWAT (early warning)} & \textbf{RCA} \\
\midrule
Question & \emph{What} / \emph{when} & \emph{Where} / \emph{why} \\
Timing   & before the fault          & after the fault \\
Output   & typed precursor alert     & causal explanation \\
Metric   & lead time, FA on drift    & localisation accuracy \\
\bottomrule
\end{tabularx}
\end{table}

\paragraph{Four operational regimes.}
A central modelling choice is that we recognise \emph{four} regimes, not three:
$\theta\in\{\theta_{\text{normal}},\theta_{\text{drift}},\theta_{\text{anomaly}},\thetadn\}$.
The fourth regime $\thetadn$ models the faulty-deployment case in which a benign drift and a real
anomaly occur simultaneously --- precisely the case where naive detectors are most dangerous.
The look-through mechanism (Stage~0) is designed to keep the signal flowing during drift rather
than zeroing it out, so that an anomaly hidden inside a drift remains observable.

\paragraph{Contributions.}
\begin{enumerate}\setlength{\itemsep}{1pt}
  \item An atomic, reproducible four-stage early-warning pipeline (Section~\ref{sec:method}),
        implemented as independent testable modules ($425{+}$ unit tests) on a real cluster.
  \item A \emph{geometric} contribution validated as H1: the STGCN latent space is structurable
        into separable fault types, silhouette $=0.782\pm0.065$ over ten seeds
        (Section~\ref{sec:results}), driven by the cosine--average clustering choice on the
        L2-normalised embedding sphere.
  \item A \emph{statistically honest} predictive headline: macro-AUROC $=0.920$
        $[0.878,0.956]$ for active-scenario typing against the independent Chaos~Mesh ground
        truth, deterministic and without the learned encoder.
  \item Evidence, via a stress-test battery (Section~\ref{sec:validity}), that the
        self-referential precursor AUROC ($0.987$) is \emph{circular}, while genuine temporal
        precursion exists against the independent target ($\Delta_{\text{far}-\text{near}}=-0.116$).
  \item Reported \emph{negative} results as contributions: look-through falsified twice
        ($p=0.27$, $p=0.37$), partial open-set recognition (OpenMax $0.55$), unstable type count.
  \item A formal OWL/RDF ontology of fault types (29 literature-anchored classes, 143
        individuals, HermiT-consistent) extracted with multivariate Transfer Entropy (KSG), never
        Granger causality.
\end{enumerate}

All numbers in this paper are traceable to the project's experiment logs; we name the dataset and
configuration for every figure, and we surface --- rather than hide --- the residual
inconsistencies we found between intermediate runs (Section~\ref{sec:limitations}).
